\section{From the selected tops to the actual arithmetic host}
\label{sec:host-tops}

\status{HC1--HC10 prove the finite, numerical and algebraic joins below.
The exact analytic inputs \texttt{ShortNorm}, \texttt{SmallOne} and
\texttt{ThinH}, and the supplied character field and pair window, remain
explicit hypotheses. In particular, constructing a host from its supplied
\texttt{HostAx} data does not construct those field data.}

\paragraph{HC1--HC2: the actual top-set choice.}
Write $n=|A|$, $o=[A:A_0]$, $n_0=|A_0|$, so $n=on_0$.
Write $Q=QS$ and $N=Qqd$ for the host moduli, and
$r=k+m-1$ for its rank parameter.
Fix $0<\varepsilon<1/52$ and $C_0\ge1$, with
$C_o=\max(C_0,2)$, before the scale or field. The actual size data give
\[
 \begin{gathered}
 \delta=x^{-\varepsilon/8},\quad h=2+8\varepsilon,
 \quad x^h\le n\le2x^h,\quad o\le C_0,\\
 k\le2x^{9/2+10\varepsilon},\quad
 x^{1/J}\le Y\le x,\quad 1\le J\le C_0.
 \end{gathered}                                      \tag{HC.1}
\]
The supplied \texttt{PairRow} says every prime
$p\le Y^{\varepsilon/5}$ splits in $\mathbb Q(\sqrt{-d})$.
An inert-class prime therefore satisfies
$p>x^{\varepsilon/(5C_o)}$. This also excludes $2$, which belongs
to the host's designated small-prime set. For $x\ge16$,
\[
 \begin{gathered}
 \log((kx^{26})^4)
 \le4\log2+(122+40\varepsilon)\log x<124\log x,\\
 L_2=2\sup_{p\,\mathrm{inert}}
       \left\lceil\log_p((kx^{26})^4)\right\rceil
 \le2\left(\frac{620C_o}{\varepsilon}+1\right).
 \end{gathered}                                      \tag{HC.2}
\]
The supremum is over natural numbers; the displayed common upper
bound proves boundedness and also covers an empty indexing set.

Apply BS5 with this last upper bound. Use the actual Frobenius steps
$h_p$ and the actual finite candidate set
\[
 P=\{p<X':p\text{ prime, inert-class},\
             \operatorname{ord}(h_p)\ge x^{\varepsilon/4}\},
 \qquad X'=\exp(x^{\varepsilon/3}).
\]
For each $t$, retain the prescribed base $b(h_p,t)$ satisfying
$b(h_p,t)+ih_p=t$ for some $i\ge0$. The badness predicate uses
the indices $0,\ldots,\operatorname{ord}(h_p)-1$ and their actual
even parity. Its base may depend on $t$. On these inert candidates
the source's \texttt{stepP} is exactly $h_p$. Thus BS5 supplies one
set $T=TT$ such that
\[
 \begin{gathered}
 \delta n/2\le |T|\le2\delta n,\\
 \bigl||T\cap(\gamma+A_0)|-\lambda_\gamma|T|\bigr|
       \le x^{-1+\varepsilon}|T|,\\
 \sum_{p\in P}\frac{\log p}{p}|B_p(T)|
       \le4C_bx^{-2}\delta n\sum_{p\in P}\frac{\log p}{p},\\
 \lambda_\gamma=
 \frac{4\mathbf1_{A_0}(\gamma)+3\mathbf1_{A\setminus A_0}(\gamma)}{3o+1}.
 \end{gathered}                                      \tag{HC.3}
\]
These are literally the size, \texttt{Frame.Shares} and
\texttt{BadOK} predicates. The source update $H[T]$ changes only
$H.TT$. Its field, \texttt{HostAx}, Frobenius law, prime classes,
$L_2$ and $P$ remain identical. Hence the data used to choose $T$
remain the data of the resulting host. The constants $C_b,x_0$
precede all these host choices.

\paragraph{HC3: transport preserves the actual bad sum.}
A point with \texttt{Pt.From} identifies its ranks, field data,
prime classes, and every Frobenius step with those of the host.
Equality of the $TT$ coordinate images, together with the
injectivity of those coordinates, identifies the actual sets.
With the canonical selector, each bad-top predicate is therefore
identical, term by term. This proves \texttt{chBad\_of\_from};
the size and share statements transfer by the same identification.
Every supplied family assembled from such points inherits
\texttt{ChosenTT} and \texttt{ChosenBad} with one threshold.
This does not supply its point or family construction.
For the same set, BS6 also gives
\[
 \sum_{p\in P}\frac{\log p}{p}|B_p(T)|
 \le16C_b|T|x^{-2+\varepsilon/3}
 \le |T|x^{-2+\varepsilon},                         \tag{HC.4}
\]
once $x^{2\varepsilon/3}\ge16C_b$.

\paragraph{HC4: what the host witness actually needs.}
A supplied \texttt{HostAx} already includes a nonempty CM-type and
an ideal $\mathfrak b$ satisfying
\[
 \mathfrak D_{E/K}=\mathfrak b\overline{\mathfrak b},
 \qquad \mathfrak b+\overline{\mathfrak b}=\mathcal O_E,
 \qquad \sigma\mathfrak b=\mathfrak b\quad
                   (\operatorname{gal}(\sigma)\in\Psi).       \tag{HC.5}
\]
Choose an embedding in that CM-type, the displayed ideal, and the
square root of $-q$ furnished by \texttt{FieldRows}. Copy them,
the supplied choices, \texttt{FieldRows}, and $q,d>4tr+1$ into
\texttt{HostNF}; its definition gives a \texttt{Host} with exactly
the original \texttt{HostAx} and choices. This proves the entire
source predicate named \texttt{admitted}. That name is not a
Lean placeholder. The substantive field facts (HC.5) remain
part of the supplied input. For the actual construction choose
$\mathfrak c_K=\mathcal O_K$: its norm is $1$, and it is primitive
because an integer $m>1$ cannot be a unit in $\mathcal O_K$.

\paragraph{HC5: the actual auxiliary primes.}
The analytic input \texttt{ShortNorm} gives constants $A,c,c'>0$
and a lower bound
\[
 \#\{p\in(X,X(1+u)]:p\text{ prime},\ \operatorname{NormOne}(d,M,p)\}
 \ge\frac{c'uX}{(4dM^2)^A\log X},                    \tag{HC.6}
\]
under its explicit hypotheses
$d\ge100$ prime, $d\equiv3\pmod4$, $M\ge3$, $X\ge1$,
$\log(4dM^2)\le c\sqrt{\log X}$ and
$e^{-c\sqrt{\log X}}\le u\le1$.
Here $\operatorname{NormOne}$ means that integers $a,b_1$ satisfy
\[
 a^2+db_1^2=4p,\qquad M\mid b_1,\qquad 2M\mid a-2-b_1.
                                                               \tag{HC.7}
\]
It is not replaced by a mere splitting condition.
If $d\mid N$ and $\log N\le3x^{\varepsilon/5}\log^2x$, the
large-$x$ consequences of (HC.6), uniformly for
$\log Z\ge8x^\varepsilon$, give at least $Z^{9/10}$ such
primes in $(Z,2Z]$. Indeed, with $v=x^{\varepsilon/5}$,
$\log(4dN^2)\le10v^2$ and $\log Z\ge8v^5$;
the denominator in (HC.6), with $u=1$, is at most
$c'Z^{1/10}$ once $x$ is large.

For the literal exceptional set
\[
 BB(r,g)=\{g\}\cup
 \{\ell:\exists\,1\le j\le r,\ \ell\text{ prime},
                      g^j<\ell^{4v_\ell(g^j-1)}\},
\]
there are at most $3r+1$ elements. Four exceptional prime-power
factors for one $j$ would already have product greater than $g^j$.
For fixed $\ell,j$, take the minimum valuation $w$ among candidates
in $[Z,2Z]$. Then $Z<(\ell^w)^4$; there are at most $2j$
roots of $z^j=1$ modulo $\ell^w$, with at most
$Z^{3/4}+1$ candidates per residue. Consequently
\[
 \begin{gathered}
 \#\{g\in[Z,2Z]:BB(r,g)\cap BB(r,g_0)\ne\varnothing\}
 \\
 \le(3r+1)[1+2r^2(Z^{3/4}+1)]\le20r^3Z^{3/4}.
 \end{gathered}                                      \tag{HC.8}
\]
Take $Z=(kb)^8$, $b=e^{x^\varepsilon}$, and $r\le x^5$.
Eventually $20r^3<Z^{3/20}$, so two of the norm primes have
disjoint $BB$ sets.

The separate analytic input \texttt{SmallOne} gives an exponent
$a_0\ge1$ and a prime $1\bmod M$ below $M^{a_0}$ for every
$M\ge2$. Apply it to $15N$, then $15Np$, then $Npp'$.
This gives distinct primes $p,p',q_0$, with
\[
 \begin{gathered}
 p,p'\equiv1\pmod{15N},\quad q_0\equiv1\pmod N,\qquad
 p,p',q_0\le N^{C_1},\\
 C_1=2a_0^3+4a_0^2+a_0.
 \end{gathered}                                      \tag{HC.9}
\]
Eventually $N^{C_1}<b<(kb)^8$, so these also differ from the
two norm primes. The five required level primes have now been
chosen with their exact separation conditions.

\paragraph{HC6: the second set is a different finite construction.}
For $L\ge2$, there is $S\subset\mathbb N$ with $4Ls<n_0$,
no relation
$j(s_2-s_1)=i(s_3-s_1)$ among distinct members with
$0<|i|,|j|\le L$ and $i\ne j$, and
\[
 |S|\ge \frac{n_0}{L}
             e^{-32\sqrt{\log n_0\log L}}.           \tag{HC.10}
\]
Here is the construction. Put $x_0=\log n_0$, $y=\log L$,
$s=\sqrt{x_0y}$, $d_0=\lfloor s/y\rfloor$,
$a=(x_0-y-3)/d_0$ and
$d=\lfloor e^a/(4L)\rfloor$ when the right side exceeds $1$;
otherwise $\{0\}$ suffices for $n_0>0$.
A sphere in $\{0,\ldots,d-1\}^{d_0}$ contains at least
$d^{d_0}/(d_0d^2)$ vectors. Encode them in base $4Ld$.
Any indicated relation has coefficients of absolute value below
the base, hence holds coordinatewise; three distinct collinear
points cannot lie on a sphere. The inequalities
$d_0\ge32$, $y+7\le a\le2s$ and
$3+d_0(7+y)+d_0+2a\le32s$ give (HC.10).

The actual $L=L_w$ satisfies
$2\le L_w\le(4084+60/\varepsilon)\log x$.
Since $n_0$ grows like $x^{2+8\varepsilon}$ up to fixed constants,
putting $E=e^{32\sqrt{\log n_0\log L_w}}$ gives, uniformly,
$E\ge2$ and $4C_oL_wE\le x^\varepsilon$ for large $x$.
The lower bound before erasure is then at least
$4x^{-\varepsilon}n\ge4$. Erase $0$; losing at most one element
gives the actual second set $T_2$ with
\[
 \begin{gathered}
 |T_2|\ge\frac{n_0}{L_w}
       e^{-64\sqrt{\log n_0\log L_w}},\qquad
 |T_2|\ge x^{-\varepsilon}n,\qquad
 4L_w|T_2|\le n_0.
 \end{gathered}                                      \tag{HC.11}
\]
The last inequality counts positive integers below $n_0/(4L_w)$.
Thus the source's absolute constant $C_{T_2}$ can be $64$.

\paragraph{HC7--HC8: the remaining elementary host joins.}
The exact rank formulas imply $r\le x^5$ and $t\le x^{400}$
for large $x$. For example, setting $y=x^{1/100}\ge2$ gives
\[
 \begin{gathered}
 f\le4y^{216},\quad k\le2y^{470},\quad r\le4y^{470},
 \quad m\ge y^{49},\quad L^2\le y^4,\\
 t\le10708y^{39967}\le y^{40000}.
 \end{gathered}
\]
The pair window gives $\log q>Y^{\varepsilon/5}/2$ and $q\mid d+1$.
These imply $q,d>4tr+1$ and $d\ge100$ once
$814\log x\le x^{\varepsilon/(5C_0)}$.
If $Q\le C_0n$, $\ell=\log Y\ge\max(2,41000C_0/\varepsilon^2)$
and $z=Y^{\varepsilon/5}$, then
\[
 \begin{gathered}
 820C_0\ell\le z,\qquad
 \log Q\le2C_0+\tfrac94\log x\le z(\ell^2-1),\\
 \log(Qqd)\le2z\ell^2.
 \end{gathered}                                      \tag{HC.12}
\]
The prescribed splitting congruences give \texttt{PairRow}, and
(HC.12) gives exactly the modulus bound needed for HC5.

Writing $b_1=Nu$, $a-2-b_1=2Nv$ in (HC.7) gives
\[
 \begin{gathered}
 \pi=\frac{a+b_1\sqrt{-d}}2
    =1+N\left(v+u\frac{1+\sqrt{-d}}2\right)\in\mathcal O_K,\\
 \pi\bar\pi=g,\qquad \pi\equiv1\pmod N,
 \qquad K=\mathbb Q(\sqrt{-d}).
 \end{gathered}                                      \tag{HC.13}
\]
Expanding the norm also gives $g\equiv1\pmod N$.
For any chosen level prime $l\equiv1\pmod N$, the supplied
Artin/Frobenius law gives trivial Frobenius and inertia. Its
decomposition group is trivial, so it splits into $[E:\mathbb Q]$
primes; quadratic reciprocity also places it in \texttt{splitK}.

The analytic \texttt{PairGood} input bounds the nontrivial twists
$\xi(\cdot/q)^i(\cdot/d)^j$ and gives their zero-free rectangles.
It yields the three \texttt{QDq} characters, the six
\texttt{Boxq} characters, and $|dlogH|\le4\log x$ directly.
For the other real characters, CRT and the unit groups at prime
powers give the precise decomposition
\[
 \begin{gathered}
 \chi=\rho(\cdot/q)^i(\cdot/d)^j,\qquad
 \operatorname{cond}(\chi)=\operatorname{cond}(\rho)q^id^j,\\
 \operatorname{cond}(\rho)\mid8Q_{\rm odd}.
 \end{gathered}                                      \tag{HC.14}
\]
Indeed a real character kills the odd-order principal units and
the squares among the dyadic units. Since
$Q\mid\operatorname{lcm}(2^{e+2},M)$ implies $Q_{\rm odd}\mid M$,
\texttt{PairGood.boxes} applies to (HC.14). Whenever $q$ or $d$
divides the conductor, $(i,j)\ne(0,0)$, giving exactly
\texttt{BoxH} and \texttt{RealBox}.

\paragraph{HC9: the exact conditional host theorem.}
Retain \texttt{ThinH}: under $Q_{\rm odd}\le C_0$, (HC.12) and
\texttt{RealBox}, every relevant primitive host character obeys
\[
 |g_\chi(y)|\le y e^{-c_6x^{\varepsilon/8}}
 \quad\left(y\ge e^{x^{\varepsilon/3}}/3\right).       \tag{HC.15}
\]
Given the three analytic inputs, choose
$C_{T_2}=64$, $C_8=\max(C_0,C_1)$, $C_5=4$, and the constants
$c_6,C_b$ before the host. Take one maximum of the preceding
thresholds. For every supplied \texttt{HostAx}, ranks and pair
window satisfying (HC.1), $Q\le C_0n$,
$0<M\le C_0$, $Q\mid\operatorname{lcm}(2^{e+2},M)$, and
\texttt{FieldRows} at $\lfloor Y^{\varepsilon/5}\rfloor$,
HC4--HC8 construct all choices and the host. Install HC2's $TT$.
The resulting host has the original \texttt{HostAx}, the exact
five-prime level set, integral norm elements with their congruences,
both second-set bounds, \texttt{QDq}, \texttt{BoxW},
$|dlogH|\le x^\varepsilon$, \texttt{ThinOK} and \texttt{BadOK}.
These are all conjuncts of
\texttt{HostGood}$(s,C_8,4,c_6,C_b,H,q_0)$.
Its shares hold separately. The expanded proof lists every
\texttt{FieldRows} assumption and checks every conjunct.

\paragraph{HC10: character windows and the ordinary specialization.}
For a primitive character $\psi$, the source field construction
supplies, for $e\ge e_0$, a host field with
$n=c2^e$, $o=\operatorname{ord}(\psi)$,
$Q\mid\operatorname{lcm}(2^{e+2},M_\psi)$,
$Q\le4M_\psi n$, and the exact kept-character, kernel and twist
data. This construction, the uniform rank constructor, the analytic
\texttt{PairEx}, and the cohomological \texttt{VWinAt}$(\psi,401)$
remain inputs here.

Their elementary window join is explicit. Put $\kappa=c^{1/h}$.
The supplied cohomological interval $[W,W^2]$, with
$W\ge\max(\kappa^3,\kappa^{400})$, gives intervals
$[Z,Z^{3/2}]$ for $Z=\kappa W$. Writing $x=\kappa x'$ gives
\[
 x^h\le c2^{e(x')}\le2x^h,\qquad
 x^{400}\le(x')^{401},\qquad W\le x'\le W^2.        \tag{HC.16}
\]
Thus the dyadic field size and the required cohomological depth
hold together. Take $C_0=J+\operatorname{ord}(\psi)+8M_\psi$ and
a threshold forcing $Y^{\varepsilon/5}\ge5+4M_\psi$ and $e\ge e_0$.
HC9 then supplies the host on every allowed scale, with one fixed
set of constants. The kernel and twist identities identify every
bounded, unkept character with a nontrivial $q,d$ twist, so
\texttt{PairGood} supplies the remaining bounds in \texttt{EAtH}.
Since $t\le x^{400}$, the supplied cohomological bound restricts
to the actual rank. This proves the literal conditional
\texttt{HostExW}$(3/2)$ endpoint of \texttt{supplyOfAdmOf}.

For $o=1$, $A_0=A$ and the shares error is zero. The source's
ordinary host-only specialization takes $M=1$, $J=1$, $Y=x$,
and uses the supplied trivial-character field and rank constructors.
It does not need \texttt{VWinAt} to obtain that host-only statement;
it also does not supply the additional cohomological bound.
The full arithmetic field and point existence questions remain
distinct from these completed selection and transport arguments.
